%% =====================================================================
%%  structures/agenda/structure.tex  --  Esqueleto de la AGENDA DIARIA
%% ---------------------------------------------------------------------
%%  Da FORMATO a los bloques de la agenda. El contenido lo compone el
%%  generador scripts/generar-agenda.py (retirado en M10 D8) o se
%%  escribe a mano a partir de templates/agenda/contenido.tex. Aqui no
%%  hay datos: solo macros de estilo, con los colores y fuentes de marca.
%%  Guardia anti-doble-carga; nombres SIN @.
%%
%%  Filosofia: se imprime RELLENADO con lo disponible y deja CASILLAS y
%%  RENGLONES en blanco para escribir a mano (no son campos de formulario).
%% =====================================================================
\ifcsname despachoEstructuraAgenda\endcsname
  \expandafter\endinput
\fi
\expandafter\def\csname despachoEstructuraAgenda\endcsname{cargada}

%% --- Casilla vacia para marcar a mano (sin amssymb) ------------------
\newcommand{\casilla}{{\color{neutral500}\framebox[10pt][c]{\rule{0pt}{7pt}}}}

%% --- Linea para escribir a mano --------------------------------------
\newcommand{\lineaescritura}{{\color{neutral500}\dotfill}}

%% --- Cabecera de la agenda: titulo + fecha + linea de control --------
\newcommand{\cabeceraagenda}{%
  \begin{center}
    {\fuentedisplaynegra\LARGE\color{guindaoficial}AGENDA DIARIA DE TRABAJO}\par
    \vspace{2pt}%
    {\fuentedatos\normalsize\color{tinta}\MakeUppercase{\ifmetadef{fecha}{\meta{fecha}}{\today}}}\par
  \end{center}
  \vspace{2pt}%
  \noindent{\color{neutral500}\rule{\textwidth}{0.6pt}}\par
  \vspace{3pt}%
  \noindent{\fuentedatos\footnotesize\color{neutral700}%
    Elaborado por:~\underline{\hspace{4.4cm}}\hfill
    Inicio:~\underline{\hspace{1.5cm}}\quad Cierre:~\underline{\hspace{1.5cm}}}\par
  \vspace{4pt}%
}

%% --- Encabezado de bloque (cuadro diferenciado: barra guinda) --------
\newcommand{\bloqueagenda}[1]{%
  \vspace{7pt}%
  \noindent\colorbox{guindaoficial}{%
    \parbox{\dimexpr\textwidth-2\fboxsep\relax}{%
      \color{white}\fuentedisplaysemi\small\strut~\MakeUppercase{#1}}}%
  \par\vspace{4pt}\noindent\ignorespaces
}

%% --- Nota breve en cursiva bajo un bloque ----------------------------
\newcommand{\notabloque}[1]{{\footnotesize\itshape\color{neutral700}#1}\par\vspace{2pt}}

%% --- Cronograma por horas (izq: hora; der: espacio para escribir) ----
\newenvironment{cronograma}{%
  \noindent\renewcommand{\arraystretch}{1.1}%
  \begin{longtable}{@{}p{1.4cm}@{}|@{\hspace{2mm}}p{\dimexpr\textwidth-2.1cm\relax}@{}}
  \hline
}{%
  \end{longtable}%
}
%%  \slot{08:00}{contenido ya impreso o vacio}
\newcommand{\slot}[2]{%
  {\fuentedatos\footnotesize\color{neutral700}#1}\rule[-0.40cm]{0pt}{0.92cm} & #2\\\hline
}
%%  evento pre-impreso dentro de un slot del cronograma
\newcommand{\slotevento}[2]{{\seriesemi\small #1}\ {\footnotesize\color{neutral700}#2}}

%% --- Caja de reunion/cita (pre-impresa + espacio de observaciones) ---
%%  \reunionbox{hora}{titulo}{lugar}{referencia}
\newcommand{\reunionbox}[4]{%
  \par\vspace{3pt}%
  \noindent\fcolorbox{neutral500}{white}{%
    \parbox{\dimexpr\textwidth-2\fboxsep-2\fboxrule\relax}{%
      \vspace{2pt}%
      \noindent{\fuentedatossemi #1}~~{\seriesemi #2}\par
      \vspace{1pt}%
      \noindent{\footnotesize\color{neutral700}Lugar: #3\hfill #4}\par
      \vspace{2pt}%
      \noindent{\footnotesize Asistió:~\casilla~Sí\quad\casilla~No\,/\,suplente
        \hfill Observaciones:~\lineaescritura}\par
      \vspace{3pt}%
    }}\par
}

%% --- Item de pendiente/documento (casilla + espacio responsable) -----
%%  \itempendiente{texto}{referencia (EXP, vence...)}
\newcommand{\itempendiente}[2]{%
  \vspace{2pt}%
  {\raggedright\noindent\casilla~~#1~{\footnotesize\color{neutral700}#2}\par}%
  \noindent\hspace{16pt}{\footnotesize Responsable:~\lineaescritura\quad Estado:~\lineaescritura}\par
}

%% --- Item de comunicacion / informativo ------------------------------
\newcommand{\itemcomunicacion}[1]{%
  \vspace{1pt}\noindent{\color{guindaoficial}\textbullet}~~#1\par
}

%% --- Renglones en blanco para escribir a mano ------------------------
%%  \rengloneslibres{6}
\newcommand{\rengloneslibres}[1]{%
  \vspace{3pt}\begingroup\color{neutral500}%
  \newcount\despachonreng \despachonreng=#1\relax
  \loop\ifnum\despachonreng>0
    \noindent\rule{\textwidth}{0.3pt}\par\vspace{0.66\baselineskip}%
    \advance\despachonreng-1
  \repeat
  \endgroup
}

%% --- Tabla de acuerdos/seguimiento (filas altas en blanco a mano) ----
\newcommand{\filaacuerdo}{ & & & \rule[-0.55cm]{0pt}{1.15cm}\\\hline}
\newcommand{\tablaacuerdos}{%
  \par\vspace{3pt}\noindent
  {\renewcommand{\arraystretch}{1.15}\setlength{\tabcolsep}{4pt}%
  \begin{longtable}{@{}|p{0.20\textwidth}|p{0.39\textwidth}|p{0.18\textwidth}|p{0.12\textwidth}|@{}}
  \hline
  {\fuentedatossemi\footnotesize~Asunto} & {\fuentedatossemi\footnotesize Acuerdo / decisión} & {\fuentedatossemi\footnotesize Responsable} & {\fuentedatossemi\footnotesize Plazo}\\\hline
  \filaacuerdo\filaacuerdo\filaacuerdo\filaacuerdo\filaacuerdo\filaacuerdo
  \end{longtable}}%
}
